\documentclass[../main]{subfiles}
\begin{document}

\chapter{Gases perfectos}
\img{images/04/gasesperfectos.jpg}{Globos inflados con Helio.}
\info{
Contenido}{
Gases perfectos,Leyes Boyle-Mariotte,Charles,Gay-Lussac
}
\actividad{Responder las siguientes preguntas y definiciones, investigar}
\begin{enumerate}
\item ¿Qué es un gas perfecto?
\item ¿Qué ley describe la relación entre la presión y el volumen de un gas en condiciones constantes de temperatura y cantidad de gas?
\item ¿Qué ley describe la relación entre el volumen y la temperatura de un gas en condiciones constantes de presión y cantidad de gas?
\item ¿Qué ley describe la relación entre la presión y la temperatura de un gas en condiciones constantes de volumen y cantidad de gas?
\item ¿Qué es la ley de Boyle-Mariotte?
\item ¿Qué es la ley de Charles?
\item ¿Qué es la ley de Gay-Lussac?
\item ¿Cómo se relacionan las tres leyes anteriores en la ecuación de estado de los gases ideales?
\item ¿Cuál es la ecuación de estado de los gases ideales?
\item ¿Qué es la constante universal de los gases y cuál es su valor en las unidades del sistema SI?
\end{enumerate}
\actividad{Resolver los siguientes ejercicios.}
\begin{enumerate}
\item Un globo de helio tiene un volumen de 5 litros a una presión de 1 atmósfera y una temperatura de 20°C. Si se le añade helio hasta que su volumen es de 10 litros a una temperatura de 30°C, ¿cuál será la nueva presión en el globo?
\item Una cámara de presión contiene 2 moles de gas nitrógeno a una temperatura de 25°C y una presión de 2 atmósferas. Si se reduce el volumen de la cámara a la mitad y se aumenta la temperatura a 50°C, ¿cuál será la nueva presión del gas?

\item Un tanque de gas propano contiene 10 litros de gas a una presión de 10 atmósferas y una temperatura de 15°C. Si se abre la válvula del tanque y se deja escapar gas hasta que la presión cae a 5 atmósferas, ¿cuál será el nuevo volumen del gas?

\item Un recipiente cerrado contiene 5 gramos de gas helio a una temperatura de 0°C y una presión de 1 atmósfera. Si se calienta el recipiente hasta que la temperatura alcanza los 27°C, ¿cuál será la nueva presión del gas?

  \item Una masa gasesosa que ocupa un volumen de $0,15m^3$ es comprimida isotérmicamente desde la presión de $1,3atm$ hasta la presión de $3,5atm$. Calcular el volumen final de la masa gaseosa.
  \item Calcular el aumento de presión que experimentará el gas encerrado en un depósito rígido, a la presión de $p=1atm$, cuando la temperatura aumenta de $17^\circ C$ a $57^\circ C$
  \item Con $1kg$ de aire a la presión atmosférica normal ($P=1,33kg/cm^2$), se observó que.
  
  $t_1=37,7^\circ C \Rightarrow V_1=0,88m^3 \quad\quad t_2=0^\circ C \Rightarrow V_2=0,7734m^3$
  
   Si el aire experimentado en el problema anterior fuera tratado como un gas ideal, ¿cuál sería la masa molecular media del gas? (ayuda: Usa la ley de los gases ideales)
  \item Calcular la presión que posee $1kg$ de oxígeno cuando ocupa un volumen $1m^3$ a la temperatura de $27^\circ C$.
  \item Calcular el volumen coupado por $m_{He}=2kg$ de helio que está comprimido a la presión de $p_{He}=5atm$ y a la temperatura de $t=-13^\circ C$.
  \item Determinar la constante del aire si a $0^\circ C$ y a la presión atmosférica normal, el volumen específico es igual a $0,7725 m^3/kg$
  \item Calcular el peso específico del amoníaco, a la presión atmosférica normal y a la temperatura de $t=-3^\circ C$
  \item El gas contenido en un cilindro ocupa el volumen de $v=0,36m^3$ a la temperatura de $t=27^\circ C$. Suponiendo que el pistón ejerce una presión constante y no existan rozamientos, calcular el volumen que ocupará la masa gaseosa si la temperatura se eleva a $t_2=72^\circ C$
  \item Calcular el peso de una masa de nitrógeno contenida en un recipiente de $v=0,6m^3$, cuando está sometida a la presión de $p=2,5atm$ y a una temperatura de $t=-13^\circ C$ 
  \item Calcular el volumen que ocupará una masa gaseosa, a la presión de $3,3atm$ y a la temperatura de $t=57^\circ C$, si a la presión de $p=5atm$ y a la temperatura de $t=-33^\circ C$ ocupa un volumen de $v=1,2m^3$.
  
  
  
\end{enumerate}
\end{document}